\documentclass[11pt]{article}
\usepackage[T1]{fontenc}
\usepackage{lmodern}
\usepackage[margin=1in]{geometry}
\usepackage{amsmath,amssymb,amsthm,mathtools}
\usepackage{booktabs,array}
\usepackage{microtype}
\usepackage{xcolor}
\usepackage[colorlinks=true,linkcolor=blue!55!black,citecolor=blue!55!black,urlcolor=blue!55!black]{hyperref}
\hypersetup{pdftitle={Sharp spectral gaps for the switch and Curveball chains on every graphical degree sequence},pdfauthor={Mark T. Short}}
\numberwithin{equation}{section}
\newtheorem{theorem}{Theorem}[section]
\newtheorem{lemma}[theorem]{Lemma}
\newtheorem{proposition}[theorem]{Proposition}
\newtheorem{corollary}[theorem]{Corollary}
\theoremstyle{definition}
\newtheorem{definition}[theorem]{Definition}
\theoremstyle{remark}
\newtheorem{remark}[theorem]{Remark}
\title{Sharp spectral gaps for the switch and Curveball chains\\on every graphical degree sequence}
\author{Mark T. Short\\[2pt]\small Independent researcher\\\small\url{https://seveneves.ai}}
\date{October 11, 2026}
\begin{document}
\maketitle
\begin{abstract}
For every graphical degree vector on $n\ge4$ labeled vertices with more than
one realization, the lazy chain that proposes switches uniformly on four
vertices has spectral gap at least $1/(6\binom n4)$, and the chain that
performs a uniform undirected Curveball trade on a uniformly chosen vertex
pair has spectral gap at least $2/\binom n2$; both constants are the worst
case over all such degree vectors, for every $n$. Consequently the
total-variation mixing times at distance $1/4$ are at most $n^6/4+1$ for the
switch chain and $n^4/16+1$ for the Curveball chain. Both gap bounds are
kernel-checked in Lean, together with a total-variation bound for the switch
chain whose prefactor is $\sqrt M/2$, where $M$ is the number of
realizations, in place of the $\sqrt{M-1}/2$ used here.
\end{abstract}
\section{Introduction}\label{sec:introduction}
Let $V=\{1,\ldots,n\}$, where $n\ge4$, and let $d=(d_1,\ldots,d_n)$ be a
graphical degree vector. Write $\Omega_d$ for the set of simple undirected
graphs on $V$ whose labeled degree vector is $d$, put $M=|\Omega_d|$, and let
$\pi$ be the uniform probability measure on $\Omega_d$. A \emph{switch}
removes two disjoint edges and inserts a different perfect matching on their
four vertices, provided both new edges are absent. A \emph{Curveball trade} at
a vertex pair $\{i,j\}$ redistributes the vertices adjacent to exactly one of
$i$ and $j$ between these two endpoints, keeping both endpoint degrees and
every other adjacency. Both moves preserve simplicity and every vertex degree.

We use the following precise discrete-time conventions. The switch chain $P$
holds with probability $1/2$. Otherwise it chooses a four-set uniformly and one
of the six ordered pairs of distinct perfect matchings on that set uniformly.
It removes the first matching and adds the second if all removed edges are
present and all added edges absent; otherwise it holds. The Curveball chain
$K$ chooses an unordered vertex pair uniformly and moves to a graph chosen
uniformly among those obtainable from the current one by a Curveball trade at
that pair, the current graph included; no additional holding probability is
imposed. Time
includes rejected proposals and holds for $P$, and counts trades for $K$. For
probability measures on $\Omega_d$ and $Q\in\{P,K\}$, set
\[
 \|\mu-\nu\|_{\mathrm{TV}}
 =\frac12\sum_{G\in\Omega_d}|\mu(G)-\nu(G)|,
 \qquad
 t_{\mathrm{mix}}^Q(\varepsilon)
 =\min\Bigl\{t\in\mathbb Z_{\ge0}:
 \max_{G\in\Omega_d}\|Q^t(G,\cdot)-\pi\|_{\mathrm{TV}}\le\varepsilon\Bigr\}.
\]
For $M>1$ write $\operatorname{gap}(Q)=1-\lambda_2(Q)$, where $\lambda_2(Q)$ is
the second largest eigenvalue of $Q$, which is symmetric
(Section~\ref{sec:setup}). Here and throughout, logarithms are
natural.

\begin{theorem}\label{thm:main}
For every graphical $d$ on $n\ge4$ vertices with $M>1$,
\begin{equation}\label{eq:gaps}
 \operatorname{gap}(P)\ge\frac{1}{6\binom n4},\qquad
 \operatorname{gap}(K)\ge\frac{2}{\binom n2}.
\end{equation}
For every $n\ge4$, both constants are best possible universal constants,
that is, the worst case over $d$: equality holds in both bounds for
$d=(2,2,1,1,0,\ldots,0)$.
\end{theorem}

Equality holds whenever $M=2$. Section~\ref{shp:section} proves this,
describes the extremal instances, and states what is not claimed for larger
realization spaces.

\begin{corollary}[Mixing times]\label{cor:mixing}
For $M>1$ and $0<\varepsilon\le1/2$, set
$A_\varepsilon=\log(1/(2\varepsilon))+\tfrac12\log(M-1)$. Then
\begin{align}
 t_{\mathrm{mix}}^P(\varepsilon)&\le\left\lceil6\binom n4 A_\varepsilon\right\rceil,
 &t_{\mathrm{mix}}^K(\varepsilon)&\le\left\lceil\frac{\binom n2}{2}
 A_\varepsilon\right\rceil,\label{eq:mixing}\\
 t_{\mathrm{mix}}^P(1/4)&\le n^6/4+1,
 &t_{\mathrm{mix}}^K(1/4)&\le n^4/16+1.\label{eq:coarse}
\end{align}
If $M=1$, both mixing times are zero.
\end{corollary}

No sharpness of these mixing bounds is asserted.

\paragraph{History and significance.}
Uniform sampling of graphs with prescribed degrees is a basic problem for
degree-constrained network models. Kannan, Tetali and Vempala~\cite{KTV99}
proved rapid mixing of the switch chain for near-regular bipartite degree
sequences and proposed the general rapid-mixing problem. Cooper, Dyer and
Greenhill~\cite{CDG07,CDG12} established polynomial mixing for regular
undirected graphs. Greenhill and Sfragara~\cite{GS18} extended this to
irregular undirected degrees under sparsity and degree conditions, and treated
directed degrees under additional hypotheses including irreducibility.
Erd\H{o}s, Greenhill, Mezei, Mikl\'os, Solt\'esz and Soukup~\cite{Unified22}
proved rapid mixing on $P$-stable families of undirected, bipartite and
directed degree sequences. These results impose restrictions on the degree
families. Tikhomirov and Youssef~\cite{TY23} proved sharp Poincar\'e bounds
for regular bipartite graphs in a specified degree range and sharp
log-Sobolev bounds for fixed degree; their setting is distinct from arbitrary
undirected degrees.

Verhelst~\cite{Verhelst08} developed trade-based MCMC sampling for binary
matrices with fixed marginals, including a uniform-stationary variant.
Carstens, Berger and Strona~\cite{CBS18} defined undirected Curveball trades
and proved convergence to the uniform distribution. Their mixing-rate
comparisons are experimental, and they left theoretical rapid mixing open.
Carstens and Kleer~\cite{CK18} used the Johnson spectrum to compare switch and
Curveball relaxation times for binary matrices, including irreducible spaces
with forbidden entries. Their comparison transfers rapid mixing from the
switch chain; it supplies no universal absolute gap, and the undirected
comparison is left open in their preprint.

The closest universal results are those of OpenAI~\cite{OAI26} and of Fu, Qin
and Wang~\cite{FQW26}. The manuscript~\cite{OAI26} is machine-checked in Lean
and has not been peer reviewed. It develops the pair-fiber framework used
here, with the triple reduction and the disjoint-cell and slice Gram methods,
and proves that $t_{\mathrm{mix}}^P(1/4)\le2n^8$ and that $P$ has spectral gap
at least $[24n^2\binom n4]^{-1}$ \cite[Theorem~1.1]{OAI26}. Fu, Qin and Wang prove sharp swap-chain and
row-pair resampling bounds for all feasible binary matrix margins; their
Theorem~5.1 also gives the three-row input used here. Arbitrary row-pair moves
do not preserve symmetric zero-diagonal adjacency matrices, so their global
matrix bound does not directly imply Theorem~\ref{thm:main}.
Theorem~\ref{thm:main} replaces the gap bound of \cite[Theorem~1.1]{OAI26} by
the best possible universal constant for the same chain, and gives the best
possible universal constant for undirected Curveball.

\paragraph{Proof strategy.}
For an unordered vertex pair $a$, let $E_a$ average a function over the
graphs reachable by a Curveball trade at $a$, put $h_a=I-E_a$ and $H=\sum_ah_a$, and let $L$ be the unnormalized
Laplacian of the switch graph on $\Omega_d$. Then $I-P=L/(12\binom n4)$ and
$K=I-H/\binom n2$ (Section~\ref{sec:setup}), so Theorem~\ref{thm:main} reduces
to $L\succeq H\succeq2(I-\Pi)$, where $\succeq$ is the order of quadratic
forms and $\Pi$ is projection onto constants. The comparison $H\preceq L$ is
the classical Johnson comparison. The second inequality follows from
$H^2\succeq2H$ and $\ker H=\mathbb R\mathbf 1$. As in
\cite[Proposition~7.1]{OAI26}, we expand
\[
 H^2=\sum_{|T|=3}H_T^2-(n-3)H+D,
\]
where $H_T$ sums $h_a$ over the three pairs $a$ inside a vertex triple $T$,
and $D$ is the
ordered sum of the products $h_ah_b$ over disjoint pairs $a,b$. OpenAI
proved $H_T^2\succeq H_T$ and $D\succeq0$, hence $H^2\succeq H$
\cite[Propositions~4.1 and~7.1, Corollary~6.3]{OAI26}. That intermediate
estimate, with the constant kernel and the classical Johnson comparison,
already gives $\operatorname{gap}(P)\ge1/(12\binom n4)$ and
$\operatorname{gap}(K)\ge1/\binom n2$, and, by the argument of
Corollary~\ref{cor:mixing}, $t_{\mathrm{mix}}^P(1/4)\le n^6/2+1$. In
particular the inverse-fourth order follows from the ingredients
of~\cite{OAI26} plus the Johnson comparison; it is not the new part of this
paper, and the improvement of Corollary~\ref{cor:mixing} over those
ingredients is the factor two.

The new step is the weighted estimate $D\succeq H$, which gives
$H^2\succeq(n-2)H-(n-3)H+H=2H$. It yields the exact factor two and the
Curveball constant $2/\binom n2$; Section~\ref{shp:section} shows separately
that both constants are attained. We prove $D\succeq H$ in two weighted
steps. For each disjoint pair, the interaction $h_ah_b+h_bh_a$ is
bounded below by twice a cross-square projection minus projections onto
centered equality indicators, each weighted by the probability $p_a$ that two
singleton labels of its fiber receive the same endpoint. On every slice, the
Johnson energy minus the weighted equality projections still dominates the
centered identity. Without the weights this slice inequality fails
(Remark~\ref{asm:weighting}). The cell analysis follows
\cite[Proposition~5.1]{OAI26} up to its final square, and the slice analysis
reuses the Gram decomposition of \cite[Lemma~6.1]{OAI26}. New here are the
weights $p_a$ with the completed square in the cell estimate, the exact
second slice-moment identity, the weighted slice inequality $G-pS\succeq h$,
and the assembly $D\succeq H$.

\paragraph{Organization.}
Section~\ref{sec:setup} defines the pair fibers and operators and identifies
both chains with them. Section~\ref{joh:section} proves the Johnson gap and
$H\preceq L$. Section~\ref{cell:section} proves the weighted inequality for
one disjoint pair, and Section~\ref{sl:section} the weighted slice
inequality. Section~\ref{asm:section} proves connectivity, $D\succeq H$ and
$H^2\succeq2H$, and tracks the constants through Theorem~\ref{thm:main} and
Corollary~\ref{cor:mixing}. Section~\ref{shp:section} proves equality
whenever $M=2$ and gives extremal instances. Section~\ref{fv:section} describes the Lean verification, and
Section~\ref{op:section} states open problems. Appendix~\ref{app:section}
proves the three-row and graph-triple inputs. All proof ingredients are
supplied in the paper.

\section{Operators and pair fibers}\label{sec:setup}
This section fixes the operator conventions and expresses both chains through
the switch Laplacian and the pair-fiber projections. All operators act on
real $L^2(\pi)$. Write $A\succeq B$ when $A-B$ is self-adjoint and has
nonnegative quadratic form. Let $\Pi f=\mathbb E_\pi f$ be projection onto
constants.

Each distinct switch neighbor of a graph has transition probability
$1/(12\binom n4)$ under $P$: the four toggled edges determine both the
four-set and the ordered pair of removed and added matchings, so this count
has no proposal multiplicity. If $L$ is the unnormalized Laplacian of the
graph of distinct switch neighbors, then
\begin{equation}\label{eq:normalization}
 Lf(G)=\sum_{G'\sim G}(f(G)-f(G')),\qquad I-P=\frac{L}{12\binom n4}.
\end{equation}

For an unordered vertex pair $a=\{i,j\}$, its \emph{fiber} fixes all edges
away from $a$, the edge $ij$, and each outside vertex's number of neighbors
in $a$. The graphs in the fiber of $G$ are exactly those obtainable from $G$
by a Curveball trade at $a$. Let $E_a$ be conditional uniform averaging on
this fiber and put
\begin{equation}\label{eq:operators}
 h_a=I-E_a,\qquad H=\sum_{a\in\binom{[n]}2}h_a,\qquad
 K=\frac{1}{\binom n2}\sum_a E_a=I-\frac{H}{\binom n2}.
\end{equation}
This is the partition-based heat-bath construction of Dyer, Greenhill and
Ullrich~\cite[Definition~1.1]{DGU14}, so $E_a$ and $h_a$ are orthogonal
projections. Thus $K$ chooses a uniform unordered pair and performs a uniform
undirected Curveball trade~\cite{CBS18}, including the possibility of
retaining the current assignment. No additional holding probability is
imposed. Both kernels are symmetric and stochastic. Also $P$ has nonnegative spectrum by
laziness, and $K$ does because it is an average of orthogonal projections.
Section~\ref{joh:section} describes each fiber as a uniform subset slice, and
Section~\ref{asm:section} combines the estimates for $H$ with
\eqref{eq:normalization} and \eqref{eq:operators}.

\section{Johnson comparison}\label{joh:section}
This section describes each pair fiber as a uniform subset slice, bounds the
pair projection $h_a$ by the exchange Laplacian of its fiber, and sums these
bounds to obtain $H\preceq L$. The gap of the exchange Laplacian is a
consequence of the classical Bernoulli--Laplace (Johnson) spectrum; see
Diaconis and Shahshahani~\cite{DS87}. We give a self-contained
local-to-global proof of the gap needed here. The identity
$\sum_aJ_a=2L$ is that of OpenAI~\cite[Lemma~7.4]{OAI26}, with the two-trade
count of Carstens, Berger and Strona~\cite[Lemma~2]{CBS18}; the exact gap
replaces the factor $n^2$ in the bound $h_a\preceq n^2J_a$
of~\cite[Lemma~7.4]{OAI26}. Section~\ref{sl:section} uses the gap on
slices, and Section~\ref{asm:section} uses $H\preceq L$.

On an $a=\{i,j\}$ fiber, let $N$ count outside vertices with exactly one
neighbor in $a$, and let $c$ count those with two. Exactly
$q=d_i-c-\mathbf 1_{ij\in E}$ singleton labels must be assigned to $i$.
Every $q$-subset of the $N$ labels is allowed, since all endpoint and outside
degrees are preserved. There is one graph per subset, so the conditional law
is uniform. Let $\tau_{kl}$ exchange the assignments of labels $k,l$, and put
\[
 J_a=\sum_{k<l}(I-\tau_{kl}).
\]
Equal assignments give zero contributions. This is the unnormalized Johnson
Laplacian, with each exchange neighbor counted once.

\begin{lemma}[Johnson gap]\label{joh:gap}
On a nontrivial uniform $q$-slice of $N$ labels, the gap of
$J=\sum_{i<j}(I-\tau_{ij})$ is $N$. Consequently $h_a\preceq J_a/2$ and
$H\preceq L$.
\end{lemma}
\begin{proof}
Write $A_a=I-\tau_{ij}$ for a label pair $a=\{i,j\}$, so $A_a/2$ is an
orthogonal projection and $A_a^2=2A_a$. Suppose first that $N\ge3$.
For each label triple $T$ set $J_T=\sum_{a\subset T,\,|a|=2}A_a$.
Fixing the subset outside $T$ leaves either one assignment, when zero or
three labels of $T$ are occupied, or three assignments, when one or two
are occupied. On the latter orbit $J_T$ is the triangle Laplacian
$3I-\mathbf 1\mathbf 1^{\mathsf T}$; on the former it is zero. Therefore
\[
 J_T^2=3J_T,\qquad \sum_{|T|=3}J_T=(N-2)J.
\]
Swaps of disjoint label pairs commute, so $A_aA_b$ is positive
semidefinite when $a\cap b=\varnothing$: it is four times the product
of two commuting orthogonal projections. Define the ordered remainder
\[
 R_J=\sum_{\substack{a,b\in\binom{[N]}2\\a\cap b=\varnothing}}A_aA_b
 \succeq0.
\]
In the sum of triple squares, each $A_a^2$ occurs $N-2$ times and each
ordered intersecting pair of distinct label pairs occurs once. Expanding
$J^2$ and using $A_a^2=2A_a$ gives
\[
 \begin{aligned}
 J^2&=\sum_{|T|=3}J_T^2-2(N-3)J+R_J\\
    &=3(N-2)J-2(N-3)J+R_J=NJ+R_J\succeq NJ.
 \end{aligned}
\]
For $N=2$ the only nontrivial slice has $q=1$, and $J^2=2J=NJ$ directly.
Thus every positive eigenvalue of the positive self-adjoint $J$ is at least
$N$. Moreover,
\[
 \langle f,Jf\rangle=\tfrac12\sum_{i<j}\lVert f-\tau_{ij}f\rVert^2.
\]
Zero energy forces invariance under every label transposition. Such
transpositions connect all $q$-subsets, so $\ker J=\mathbb R\mathbf 1$.
Finally, with $x_i$ the membership indicator of label $i$, direct summation
gives $Jx_i=\sum_{j\ne i}(x_i-x_j)=Nx_i-q$. Since $1\le q<N$,
$x_i-q/N$ is a nonzero centered eigenfunction with eigenvalue $N$.
The gap is therefore exactly $N$.

On a nontrivial fiber $N\ge2$, so $h_a\preceq J_a/N\preceq J_a/2$;
on one-state fibers both operators vanish. Every switch is realized by
exactly two singleton exchanges: their indexing pairs are the two parts of
the unique bipartition of its toggled four-cycle. For each indexing pair the
outside labels are determined. Conversely, exchanging a label $k$ assigned
to $i$ with a label $l$ assigned to $j$ toggles exactly the four edges
$ik,jl,il,jk$, which is a valid switch, so every nontrivial singleton
exchange is a switch. Thus, by off-diagonal entries and zero row sums,
\begin{equation}\label{joh:multiplicity}
 \sum_a J_a=2L,\qquad H\preceq\tfrac12\sum_aJ_a=L.
\end{equation}
\end{proof}

\begin{remark}
The Johnson proof has the same local-to-global shape as the argument for
$H^2\succeq2H$: a triple square identity or inequality, a regrouping of
interactions, and a constant-kernel argument. Disjoint label swaps commute,
making $R_J$ positive directly. Disjoint graph-pair resamplings need not
commute; Sections~\ref{cell:section}--\ref{asm:section} control their
interaction by the weighted estimates.
\end{remark}

\section{Weighted cell inequality}\label{cell:section}
This section bounds the interaction $h_ah_b+h_bh_a$ of two disjoint pair
resamplings from below. The proof follows the disjoint-pair proposition of
OpenAI~\cite[Proposition~5.1]{OAI26} except for the final weighted square.
The new ingredient is the fiber-constant weight $p_a$ on the error terms,
which retains the identity $\alpha\beta=\rho^2$ in that square. Section~\ref{asm:section} sums the lemma over disjoint pairs
and pays for the error terms with Lemma~\ref{sl:weighted}.

On an $a$-fiber with parameters $(N,q)$ define the fiber-constant multiplier
\begin{equation}\label{cell:p}
 p_a=\frac{q(q-1)+(N-q)(N-q-1)}{N(N-1)}\quad(N\ge2),
 \qquad p_a=0\quad(N\le1).
\end{equation}
It is the probability that two distinct singleton labels have equal
assignments. For disjoint $a,b$, let $T_{a,b}$ toggle the cross square
between $a$ and $b$ when it is alternating, and fix the graph otherwise. It is
a measure-preserving involution. Set $g_{a,b}=(I-T_{a,b})/2=g_{b,a}$, an
orthogonal projection. Let $R_{a,b}$ act on each $a$-fiber as projection onto
the centered indicator that the two labels of $b$ have equal assignments; it
is zero unless both are singletons, and also zero if that centered indicator
vanishes. Consequently $p_aR_{a,b}$ is self-adjoint and positive
semidefinite.

\begin{lemma}[Weighted disjoint-pair interaction]\label{cell:interaction}
For disjoint vertex pairs $a,b$,
\begin{equation}\label{cell:bound}
 h_ah_b+h_bh_a\succeq2g_{a,b}-p_aR_{a,b}-p_bR_{b,a}.
\end{equation}
\end{lemma}
\begin{proof}
Write $a=\{i,j\}$, $b=\{k,l\}$. Partition graphs into cells by fixing the
graph outside these four vertices, the edges $ij,kl$, and each outside
vertex's separate neighbor counts in $a$ and $b$. Both updates preserve these
data. Let $A_0,B_0$ be the two outside singleton sets. After subtracting fixed
double-neighbor and internal-edge contributions, let $q_0,q_1$ be the residual
degrees of $i,k$. A cell is exactly the equally likely encodings
$(F,U,W)$, where $F$ is a binary $2\times2$ cross matrix of fixed total $s$,
$r,c$ are its first row and first column sums, and
\[
 U\subseteq A_0,\quad |U|=q_0-r,\qquad
 W\subseteq B_0,\quad |W|=q_1-c.
\]
The total $s$ is fixed by the sum of the degrees in $a$. The displayed
conditions enforce the first endpoint degrees; the second endpoint degrees
then follow from the totals. All outside degrees are fixed as well.
Even if $A_0,B_0$ overlap, their assignments use distinct edges and impose no
additional constraint. This proves the claimed bijection and uniform law.

Within a cell the full $a$-fiber signature is $X=(c,W)$ and the full
$b$-fiber signature is $Y=(r,U)$. Let $\mathcal A,\mathcal B$ be their
function spaces, including constants, and
$\mathcal J=(\mathcal A+\mathcal B)^\perp$. The cross involution preserves
$X,Y$, so $\operatorname{ran}g_{a,b}\subseteq\mathcal J$ and
$P_{\mathcal J}\succeq g_{a,b}$.

\paragraph{Product cases.}
For $s=0,4$ the cross matrix is fixed. For $s=1,3$, its unique one or unique
zero is determined by independent row and column choices and the subset
weights factor. Thus $X,Y$ have a product law in these cases. For $s=2$ the
cross-matrix multiplicities are
\begin{equation}\label{cell:crossmatrix}
 (m_{rc})_{r,c=0}^2=\begin{pmatrix}0&1&0\\1&2&1\\0&1&0\end{pmatrix}.
\end{equation}
Set $\Delta_r=\{r\ne1\}$ and $\Delta_c=\{c\ne1\}$; they are disjoint. If either event
is absent, one total is fixed and the remaining weights factor, again giving
a product law. In every product case the conditional projections commute,
their product is projection onto constants, and $h_ah_b=P_{\mathcal J}$. This
proves \eqref{cell:bound} there.

\paragraph{The rank-one case.}
In the remaining case put $p_r=\Pr(\Delta_r)>0$ and $p_c=\Pr(\Delta_c)>0$.
The occurrence of $\Delta_c$ ensures a feasible $U$ of size $q_0-1$, and that of
$\Delta_r$ ensures a feasible $W$ of size $q_1-1$. Hence the central cross matrices
are feasible, and $p_r+p_c<1$. Define
\[
 u=\frac{\mathbf 1_{\Delta_c}-p_c}{\sqrt{p_c(1-p_c)}},\quad
 v=\frac{\mathbf 1_{\Delta_r}-p_r}{\sqrt{p_r(1-p_r)}},\quad
 \langle u,v\rangle=-\rho=-\sqrt{\frac{p_cp_r}{(1-p_c)(1-p_r)}}.
\]
In particular $0<\rho<1$. Conditional on $X=(c,W)$, the weight of each
feasible $(r,U)$ is proportional to $m_{rc}$, with normalizer
$\sum_t\binom{|A_0|}{q_0-t}m_{tc}$. By \eqref{cell:crossmatrix} this
conditional law depends only on whether $c=1$; the reverse statement holds
for $r$. The centered correlation therefore has rank one:
$E_av=-\rho u$, $E_bu=-\rho v$. The centered parts of $\mathcal A$ orthogonal
to $u$, and of $\mathcal B$ orthogonal to $v$, are orthogonal to the opposite
space. They, constants, $\operatorname{span}(u,v)$, and $\mathcal J$ form an
orthogonal invariant decomposition.

The interaction is zero on constants and those remaining row spaces, and is
$2I$ on $\mathcal J$. On the plane write the component of a general $f$ as
$xu+yv$. Its interaction quadratic form is
\[
 \Phi(f)=2\langle h_a(xu+yv),h_b(xu+yv)\rangle
       =2xy\rho(1-\rho^2),
\]
and its defect images are $y(v+\rho u)$ and $x(u+\rho v)$, with squared norms
$y^2(1-\rho^2)$ and $x^2(1-\rho^2)$.

\paragraph{The weighted square.}
The line spanned by $v+\rho u$ lies in $\operatorname{ran}R_{a,b}$: it
vanishes for $c\ne1$, and on each $c=1$ fiber it is a multiple of the centered
equal-assignment indicator. The analogous assertion holds for $u+\rho v$ and
$R_{b,a}$. Because these lines lie in the plane, orthogonal components of the
full $f$ do not affect its projection onto them. Therefore
\[
 \lVert R_{a,b}f\rVert^2\ge y^2(1-\rho^2),\qquad
 \lVert R_{b,a}f\rVert^2\ge x^2(1-\rho^2).
\]
Every fiber supporting $R_{a,b}$ has $(N,q)=(|A_0|+2,q_0)$. Hence $p_a$
takes the value
\[
 \alpha:=\Pr(\Delta_r\mid c=1)=\frac{p_r}{1-p_c}
\]
on every $a$-fiber where $R_{a,b}$ is nonzero, and $p_b$ takes the value
$\beta:=p_c/(1-p_r)$ on every $b$-fiber where $R_{b,a}$ is nonzero.
Indeed $c=1$ makes both labels of $b$ singletons in the $a$-fiber, and they
have equal assignments exactly when $r\ne1$. The uniform slice probability
is precisely \eqref{cell:p}, independently of $W$. Now $\alpha\beta=\rho^2$,
so
\begin{align*}
 \Phi(f)+\langle f,p_aR_{a,b}f\rangle+\langle f,p_bR_{b,a}f\rangle
 &\ge(1-\rho^2)(2\rho xy+\alpha y^2+\beta x^2)\\
 &=(1-\rho^2)(\sqrt\beta x+\sqrt\alpha y)^2\ge0.
\end{align*}
The remaining contribution $2\lVert P_{\mathcal J}f\rVert^2$ is at least
$2\lVert g_{a,b}f\rVert^2$. This proves the inequality on every cell and
hence globally.
\end{proof}

\section{Weighted slice inequality}\label{sl:section}
This section concerns a uniform subset slice and has no graph interpretation.
It compares the Johnson energy with the projections onto centered equality
indicators, weighted by the equality probability $p$ of \eqref{cell:p}. The
proof reuses the Gram decomposition and the first slice-moment computation
of OpenAI~\cite[Lemma~6.1]{OAI26}, which rest on Filmus's decomposition of
the Johnson scheme~\cite[Lemmas~18 and~24]{Filmus16}. The second
slice-moment identity and the weighted inequality are new; without the weight
$p$ the inequality fails (Remark~\ref{asm:weighting}).
Section~\ref{asm:section} applies the lemma on every pair fiber.

On a uniform $q$-slice of $N$ labels let $x_i$ be the membership indicators,
$h=I-\Pi$, where here $\Pi$ denotes the slice average, $G=J/2$, and let $p$
be given by \eqref{cell:p}. Set
\[
 z_{ij}=\mathbf 1_{x_i=x_j}-p,\qquad
 S=\sum_{i<j}Q_{ij},
\]
where $Q_{ij}$ is projection onto $z_{ij}$, interpreted as zero if
$z_{ij}=0$.

\begin{lemma}[Weighted slice inequality]\label{sl:weighted}
On every uniform slice, $G-pS\succeq h$.
\end{lemma}
\begin{proof}
One-state slices have $G=S=h=0$. For $N=2,q=1$, $p=S=0$ and $G=h$.
Otherwise $N\ge3$, $1\le q<N$, and $0<p<1$. Put $r=N-q$ and $m_N=\binom N2$.
Map the Euclidean space of pair coefficients into the slice space by
\[
 Bc=\sum_{i<j}c_{ij}\frac{z_{ij}}{\sqrt{p(1-p)}}.
\]
Then $S=BB^*$. The coefficient space decomposes orthogonally as
\[
 C_0=\mathbb R\mathbf 1,\quad C_1=\{c_{ij}=u_i+u_j:\ \sum_i u_i=0\},\quad
 C_2=\{c:\ \sum_{j\ne i}c_{ij}=0\text{ for every }i\}.
\]
Here $c_{ji}=c_{ij}$. The map $u\mapsto(u_i+u_j)$ is injective for $N\ge3$;
the orthogonal complement of its whole image is exactly $C_2$.
The dimensions are $1,N-1,m_N-N$, so $C_2=0$ if $N=3$.

\paragraph{The Gram scalars.}
The number of equal pairs is fixed; hence $BC_0=0$. By label symmetry the
entries of $B^*B$ depend only on the intersection size of the indexing pairs.
For a fixed pair $ij$, the sums of $c_{kl}$ over other pairs meeting it and
over disjoint pairs are respectively
$(N-4)c_{ij},-(N-3)c_{ij}$ on $C_1$, and $-2c_{ij},c_{ij}$ on $C_2$.
These follow by summing the row sums (which are $(N-2)u_i$ on $C_1$)
and using the zero total sum. Thus $B^*B$ acts as scalars $\mu_1,\mu_2$.
Its diagonal entries are one. The identities
\begin{align*}
 \sum_{i<j}(u_i+u_j)z_{ij}&=(2q-N)\sum_i u_ix_i,\\
 \operatorname{Var}\Big(\sum_i u_ix_i\Big)&=\frac{qr}{N(N-1)}\sum_i u_i^2,\qquad
 \sum_{i<j}(u_i+u_j)^2=(N-2)\sum_i u_i^2
\end{align*}
follow from $\sum_i u_i=0$, $\sum_i x_i=q$, and
$\mathbb E x_i=q/N$, $\mathbb E x_ix_j=q(q-1)/(N(N-1))$. Since
$1-p=2qr/(N(N-1))$, they give the first formula below; the trace identity
$(N-1)\mu_1+(m_N-N)\mu_2=m_N$ gives the second:
\[
 p\mu_1=\frac{(N-2q)^2}{2(N-2)},\qquad
 p\mu_2=\frac{2(q-1)(r-1)}{(N-2)(N-3)}\quad(N\ge4).
\]

\paragraph{Comparison on the images.}
The mutually orthogonal images $U_i=BC_i$ exhaust $\operatorname{ran}S$, and
$S$ acts there by $\mu_i$. Direct transposition summation gives
\[
 Gx_i=(Nx_i-q)/2,\qquad
 G(x_ix_j)=(N-1)x_ix_j-(q-1)(x_i+x_j)/2.
\]
The image $U_1$ is linear. For $C_2$, expansion of
$z_{ij}=1-p-x_i-x_j+2x_ix_j$ cancels the constant and linear coefficients;
its row sums also cancel the lower terms in the second displayed formula.
Consequently $G$ acts by $N/2$ on $U_1$ and by $N-1$ on $U_2$.
Zero images require no eigenvalue assertion. On the nonzero images the
eigenvalues of $G-pS-h$ are nonnegative: on $U_1$,
\[
 \frac N2-p\mu_1-1=\frac{2(q-1)(r-1)}{N-2}\ge0;
\]
on $U_2$, using $(q-1)(r-1)\le(N-2)^2/4$,
\[
 0\le p\mu_2\le\frac{N-2}{2(N-3)}\le1,\qquad
 N-1-p\mu_2-1\ge N-3>0.
\]
On the orthogonal complement of constants and $U_1\oplus U_2$, $S=0$, while
Lemma~\ref{joh:gap} gives $G\succeq(N/2)I\succeq I$. These subspaces
are invariant under the self-adjoint $G$; constants are killed by all the
operators. This completes the proof, including all boundary cases.
\end{proof}

\section{Assembly}\label{asm:section}
This section sums the weighted cell and slice inequalities into $D\succeq H$
and $H^2\succeq2H$, identifies the kernel of $H$, and tracks the constants
through Theorem~\ref{thm:main} and Corollary~\ref{cor:mixing}. The triple
estimate is Lemma~\ref{app:triple}, proved in Appendix~\ref{app:section}.

\subsection{The global operator inequality}
For a vertex triple $T$ let $H_T=\sum_{a\subset T,\,|a|=2}h_a$, and define
\[
 D=\sum_{\substack{a,b\in\binom{[n]}2\\a\cap b=\varnothing}}h_ah_b.
\]
This is an ordered sum and is self-adjoint; no commutation is assumed.
Using $h_a^2=h_a$, each diagonal term occurs $n-2$ times in the triple sum,
and each ordered pair of distinct intersecting pairs occurs once. Thus we
obtain the expansion of OpenAI~\cite[Proposition~7.1]{OAI26}, which follows
Caputo's squared-generator method~\cite[Section~1.1,
Equations~(1.13)--(1.15)]{Caputo08}:
\begin{equation}\label{asm:expansion}
 H^2=\sum_{|T|=3}H_T^2-(n-3)H+D,\qquad
 \sum_{|T|=3}H_T=(n-2)H.
\end{equation}

\begin{proposition}[Disjoint and global estimates]\label{asm:global}
The disjoint-pair sum satisfies $D\succeq H$, and consequently
$H^2\succeq2H$.
\end{proposition}
\begin{proof}
On an $a$-fiber, pairs $b$ containing a nonsingleton label contribute zero to
both $g_{a,b}$ and $R_{a,b}$. The other pairs index exactly the
singleton-label transpositions and equality lines. Therefore
\[
 G_a:=\sum_{b\cap a=\varnothing}g_{a,b}=J_a/2,\qquad
 S_a:=\sum_{b\cap a=\varnothing}R_{a,b}=S\text{ on that fiber}.
\]
Sum Lemma~\ref{cell:interaction} over \emph{unordered} disjoint $\{a,b\}$.
Each such pair contributes two copies of $g_{a,b}$, giving exactly
$\sum_a G_a$, and each directed error occurs once. Lemma~\ref{sl:weighted},
averaged over fibers, now gives
\[
 D\succeq\sum_a(G_a-p_aS_a)\succeq\sum_a h_a=H.
\]
Using the triple estimate of Lemma~\ref{app:triple} in
\eqref{asm:expansion} yields
\[
 H^2\succeq(n-2)H-(n-3)H+D\succeq2H.
\]
\end{proof}

\begin{remark}\label{asm:weighting}
The weights cannot be removed from Lemma~\ref{sl:weighted}: on $q=1$ slices
with $N\ge3$, $G=S$, so $G-S\succeq h$ is false, whereas $G-pS=h$. Retaining
$\alpha\beta=\rho^2$ in the cell square leaves the slice energy needed for
the strengthened global bound.
\end{remark}

\subsection{Connectivity and the kernel}
The following lemma is the connectivity lemma and kernel corollary of
OpenAI~\cite[Lemma~7.2, Corollary~7.3]{OAI26}, whose connectivity argument
is Havel's switching argument~\cite[Theorem~3]{Havel55}. We include the
proof.

\begin{lemma}[Connectivity and kernel]\label{asm:kernel}
Switches connect $\Omega_d$, and $\ker H=\mathbb R\mathbf 1$.
\end{lemma}
\begin{proof}
Let $R$ be the current set of unfrozen vertices, initially $V$, and write
$d^R_w$ for the degree of $w$ in the graph induced on $R$ and $N_R(w)$ for
its neighborhood there. Switches inside $R$ leave every residual degree
$d^R_w$ unchanged, and hence also the prescription below. Choose the first
vertex $v$ of $R$ by label and prescribe its neighborhood in $R$ to be the
$d^R_v$ vertices of $R\setminus\{v\}$ of largest residual degree, breaking
ties by label. If $i$ is an undesired current neighbor and $j$ a prescribed
missing neighbor, then $d^R_j\ge d^R_i$. The sets
$A=N_R(j)\setminus\{i\}$ and $B=N_R(i)\setminus\{j\}$ satisfy $|A|\ge|B|$
and $v\in B\setminus A$. Hence some $k\in A\setminus B$ exists. These four
vertices are distinct and $\{vi,jk\}\to\{vj,ik\}$ is a valid switch, gaining
a prescribed neighbor without losing another. Repeat, then freeze $v$ and
recurse. The resulting residual degree vector depends only on the preceding
vector and the fixed prescription. Every realization therefore reaches the
same graph. With fewer than four unfrozen vertices a discrepancy would still
force four distinct vertices, so none occurs.

Since the $h_a$ are orthogonal projections,
$\langle f,Hf\rangle=\sum_a\lVert h_af\rVert^2$. Zero energy forces $f$ to be
constant on every fiber and hence along every switch, by the multiplicity
count behind \eqref{joh:multiplicity}. Connectivity proves the kernel
statement, and constants have zero energy because $h_a\mathbf 1=0$.
\end{proof}

\begin{corollary}[Operator comparison]\label{asm:comparison}
The operators satisfy $L\succeq H\succeq2(I-\Pi)$.
\end{corollary}
\begin{proof}
By Proposition~\ref{asm:global}, the positive self-adjoint $H$ has each
eigenvalue either zero or at least two. Lemmas~\ref{asm:kernel}
and~\ref{joh:gap} imply $L\succeq H\succeq2(I-\Pi)$.
\end{proof}

\subsection{The prescribed chains}
\begin{proof}[Proof of Theorem~\ref{thm:main}]
By Corollary~\ref{asm:comparison} and \eqref{eq:normalization},
\[
 I-P=\frac{L}{12\binom n4}\succeq\frac{2(I-\Pi)}{12\binom n4}
 =\frac{I-\Pi}{6\binom n4},
\]
and by \eqref{eq:operators}, $I-K=H/\binom n2\succeq2(I-\Pi)/\binom n2$. Both
kernels are symmetric and fix constants, so these are the two bounds in
\eqref{eq:gaps}. By Proposition~\ref{shp:extremal},
$d=(2,2,1,1,0,\ldots,0)$ has $M=2$, so Proposition~\ref{shp:two} gives
equality in both bounds.
\end{proof}

\begin{proof}[Proof of Corollary~\ref{cor:mixing}]
The case $M=1$ is immediate: the unique graph carries all of $\pi$. For
$M>1$ and either kernel, the centered density of a point mass has squared
norm $M-1$. Reversibility and the nonnegative spectrum give, for a valid gap
lower bound $\gamma_0$ and every starting graph,
\begin{equation}\label{asm:tv}
 \|Q^t(G,\cdot)-\pi\|_{\mathrm{TV}}
 \le\tfrac12\sqrt{M-1}(1-\gamma_0)^t
 \le\tfrac12\sqrt{M-1}e^{-\gamma_0t}.
\end{equation}
Substitute $\gamma_0=1/(6\binom n4)$ or $2/\binom n2$ to obtain
\eqref{eq:mixing}. Since $M\le2^{\binom n2}$,
\[
 A_{1/4}\le\left(1+\frac{n(n-1)}4\right)\log2\le n^2,
 \qquad 6\binom n4\le\frac{n^4}{4}.
\]
For the second coarse constant, $n\ge4$ gives
\[
 \frac{\binom n2}{2}A_{1/4}
 \le\left(\frac{n^2}{4}+\frac{n^4}{16}\right)\log2
 =\frac{n^4}{16}(1+4/n^2)\log2\le\frac{n^4}{16}.
\]
Here $(1+4/n^2)\log2\le(5/4)\log2<1$, because
$\log2=\int_1^2 dx/x<3/4$ by the trapezoidal bound for the strictly convex
function $1/x$. Finally $\lceil x\rceil\le x+1$ proves \eqref{eq:coarse}.
\end{proof}

\section{Best possible constants}\label{shp:section}
This section proves that the constants of Theorem~\ref{thm:main} are best
possible universal constants for every $n\ge4$, that is, the worst case over
$d$. Equality holds for every degree vector with exactly two realizations,
and such vectors exist for every $n\ge4$.

\begin{proposition}[Two realizations]\label{shp:two}
If $M=2$, equality holds in both bounds of \eqref{eq:gaps}.
\end{proposition}
\begin{proof}
By Lemma~\ref{asm:kernel} the two realizations are switch neighbors, so $L$
is the Laplacian of one edge, with gap two. Within a pair fiber each
reassigned label toggles two edges, so two graphs of one fiber that differ in
four edges differ by one singleton exchange. Hence a pair fiber contains both
realizations only when the pair is one of the two parts of the bipartition of
the toggled four-cycle, by the multiplicity count in the proof of
Lemma~\ref{joh:gap}. Every other pair fiber has one state, so exactly two
$h_a$ are nonzero, and each equals $I-\Pi$ on the two-state space. Thus
$H=2(I-\Pi)$, and \eqref{eq:normalization}--\eqref{eq:operators} give
$I-P=(I-\Pi)/(6\binom n4)$ and $I-K=2(I-\Pi)/\binom n2$, which is equality in
both parts of \eqref{eq:gaps}.
\end{proof}

\begin{proposition}[Extremal instances]\label{shp:extremal}
For every $n\ge4$, the degree vector $d=(2,2,1,1,0,\ldots,0)$ has $M=2$. The
degree vector $(3,3,2,1,1)$ has $M=2$ and neither an isolated nor a
dominating vertex.
\end{proposition}
\begin{proof}
Take $d=(2,2,1,1,0,\ldots,0)$. The edge $12$ must be present: otherwise both
vertices $1,2$ would require both $3,4$, contradicting their degrees. The only
realizations are
\[
 \{12,13,24\},\qquad \{12,14,23\}.
\]
They differ by one switch, and the two nontrivial pair fibers are indexed by
$\{1,2\}$ and $\{3,4\}$; each is the whole two-state space.

Take $d=(3,3,2,1,1)$. The edge $12$ must be present: otherwise both vertices
$1,2$ would be adjacent to all of $3,4,5$, giving vertex $4$ degree two. The
four remaining edges at $1$ and $2$ end in $\{3,4,5\}$, whose degree sum is
four, so no edge lies inside $\{3,4,5\}$ and vertex $3$ is adjacent to both
$1$ and $2$. Each of $1,2$ then has one neighbor in $\{4,5\}$, and these
neighbors are distinct because vertices $4,5$ have degree one. The only
realizations are therefore
\[
 \{12,13,23,14,25\},\qquad \{12,13,23,15,24\}.
\]
No degree is zero, and none equals $n-1=4$.
\end{proof}

The instance $(2,2,1,1,0,\ldots,0)$ is a two-element realization space with
$n-4$ isolated vertices appended. Nothing is claimed here about the best
constants over degree vectors whose realization spaces are larger; whether
equality or near-equality occurs for large $M$ is open
(Section~\ref{op:section}). For regular bipartite graphs in a specified
degree range, Tikhomirov and Youssef~\cite{TY23} proved sharp Poincar\'e
bounds for the switch chain, and sharp log-Sobolev bounds for fixed degree.

\section{Formal verification}\label{fv:section}
The two gap lower bounds of Theorem~\ref{thm:main} and a weaker
total-variation bound for $P$ are checked in Lean~4 with Mathlib. The switch
statements use the definitions of the Lean formalization accompanying
\cite{OAI26} at commit \texttt{adc7f1241b42e322a6451854ab7e4b4c146bf78a},
unchanged. The toolchain pins are as follows.
\begin{center}
\begin{tabular}{ll}
\toprule
Lean & \texttt{leanprover/lean4:v4.34.1}\\
Mathlib & \texttt{d13f23b723b8a846827a245b89c10fc7d3f11612}\\
\bottomrule
\end{tabular}
\end{center}

\paragraph{Switch contracts.}
The file \texttt{Statement.lean} imports only the pinned definitions and
states two contracts. Both assume $n\ge4$, a degree vector $d$ with
$d_v\le n-1$ for every $v$, $d$ graphical, and $M>1$, where $\Omega_d$ is the
pinned state space \texttt{GraphState}~$n$~$d$ and $M$ is
\texttt{stateCount}~$n$~$d$. Let $P$ be the pinned \texttt{switchKernel},
$\operatorname{Var}_\pi$ the pinned \texttt{uniformVariance}, and
$\mathcal E_P(f)=\langle f,(I-P)f\rangle$ the pinned \texttt{dirichletEnergy}.
The contracts read
\begin{gather*}
 \frac{1}{6\binom n4}\operatorname{Var}_\pi(f)\le\mathcal E_P(f)
 \quad\text{for every }f\colon\Omega_d\to\mathbb R,\\
 \|P^t(G,\cdot)-\pi\|_{\mathrm{TV}}
 \le\frac{\sqrt M}{2}\exp\Bigl(-\frac{t}{6\binom n4}\Bigr)
 \quad\text{for every }t\in\mathbb Z_{\ge0}\text{ and }G\in\Omega_d.
\end{gather*}
The pinned gap predicate is this Poincar\'e form: $\gamma$ times the
variance is at most the Dirichlet energy for every $f$. For a reversible
kernel the largest such $\gamma$ is $1-\lambda_2$, so the first contract is
the switch bound in \eqref{eq:gaps}. The second contract has prefactor
$\sqrt M/2$, weaker than the $\sqrt{M-1}/2$ of \eqref{asm:tv}. The theorems
\texttt{SharpSwitchGap.\allowbreak switch\_chain\_gap} and
\texttt{SharpSwitchGap.\allowbreak switch\_chain\_tv\_bound} prove the two
statements, the latter by the general gap-to-total-variation argument of the
upstream formalization. The theorems
\texttt{SharpSwitchGap.\allowbreak Baseline.\allowbreak sharp\_gap} and
\texttt{SharpSwitchGap.\allowbreak Baseline.\allowbreak sharp\_tv\_bound}
check these proofs against the two contracts. A synchronization script verifies the Git blob hash of
every copied upstream file and matches the definitions and upstream theorem
types against the pinned upstream statements.

\paragraph{Curveball contract.}
The project defines the Curveball kernel \texttt{Curveball.kernel} on the
pinned state space: it chooses a uniform unordered pair and then a
uniform state among the pinned \texttt{fiberStates} of that pair, the current
state included. Both identities in \eqref{eq:operators} are proved by the
theorems \texttt{kernel\_apply} and \texttt{operator\_average};
further theorems prove that the kernel is nonnegative, stochastic and
symmetric, and that its operator is positive semidefinite. With
$\mathcal E_K(f)=\langle f,(I-K)f\rangle$, the theorem \texttt{Curveball.gap}
proves, for $n\ge4$ and graphical $d$,
\[
 \frac{2}{\binom n2}\operatorname{Var}_\pi(f)\le\mathcal E_K(f)
 \quad\text{for every }f\colon\Omega_d\to\mathbb R.
\]
The Curveball kernel and its energy and gap predicates are defined in this
project on the pinned graph space rather than taken from upstream, so this
contract rests on a larger trusted base than the switch contracts.

\paragraph{Checked argument.}
Table~\ref{fv:declarations} lists the intermediate results of the checked
argument. The triple estimate $H_T^2\succeq H_T$ enters as the pinned
upstream formal result; the proofs in Appendix~\ref{app:section} are not
separately formalized. The audited endpoints depend only on the axioms
\texttt{propext}, \texttt{Classical.choice} and \texttt{Quot.sound}.

\begin{table}[ht]
\centering
\small
\begin{tabular}{ll}
\toprule
Statement & Lean declaration\\
\midrule
$H\preceq L$ (Lemma~\ref{joh:gap}) &
\texttt{SharpSwitchGap.SharpComparison.generator\_le\_laplacian}\\
Lemma~\ref{sl:weighted} &
\texttt{SharpSwitchGap.WeightedSlice.weighted\_centered}\\
Lemma~\ref{cell:interaction}, summed over cells &
\texttt{SharpSwitchGap.WeightedCell.actual\_named\_interaction}\\
$D\succeq H$ &
\texttt{SharpSwitchGap.SharpAssembly.disjoint\_ge\_generator}\\
$H^2\succeq2H$ &
\texttt{SharpSwitchGap.SharpAssembly.generator\_square}\\
$H\succeq2(I-\Pi)$ &
\texttt{SharpSwitchGap.SharpAssembly.centered\_poincare}\\
\bottomrule
\end{tabular}
\caption{Intermediate results and the Lean declarations that check them.}
\label{fv:declarations}
\end{table}

\paragraph{Reproduction.}
With the Lean toolchain available, run the following from the \texttt{lean}
directory of the repository at \url{https://github.com/sevenevesai/SharpSwitchGap}.
\begin{verbatim}
python verify_vendor.py
lake exe cache get
lake build
lake env lean SharpSwitchGap/B5Checks.lean
lake env lean SharpSwitchGap/B4Checks.lean
\end{verbatim}
The first command checks the vendored upstream files against the Git blob
hashes of the pinned commit. The two audits print the axiom dependencies of
the Curveball and switch endpoints and of the two contract checks. Success
means exit code zero, no \texttt{sorry} diagnostic, and only the three axioms
above in every printed dependency list.

\paragraph{Not formalized.}
The following results of the paper are not formalized:
\begin{itemize}
\item the equality cases, Propositions~\ref{shp:two} and~\ref{shp:extremal};
\item the exact Johnson gap equality in Lemma~\ref{joh:gap}; the lower bound
 is checked;
\item the $\sqrt{M-1}/2$ prefactor in \eqref{asm:tv}; the checked switch
 contract uses $\sqrt M/2$ with the same gap constant;
\item a total-variation estimate for the Curveball chain $K$;
\item the mixing-time bounds of Corollary~\ref{cor:mixing}.
\end{itemize}

\section{Open problems}\label{op:section}
In the normalizations of Section~\ref{sec:introduction}, the worst-case
relaxation times have orders $n^4$ for $P$ and $n^2$ for $K$, by
Theorem~\ref{thm:main}. At fixed error, Corollary~\ref{cor:mixing} gives only
$O(n^6)$ and $O(n^4)$ mixing-time bounds, respectively. The loss comes from
the $\log M$ factor in the gap-to-total-variation estimate. A log-Sobolev or
entropy-contraction analogue of the pair-fiber argument strong enough to
avoid this factor is not proved here.

Directed degree sequences require a separate extension. A formulation must
specify an irreducible chain, since ordinary directed switches need not
connect every realization space; the directed results
in~\cite{GS18,Unified22,CK18} have their own hypotheses. The weighted
interaction argument is not extended here to the forbidden diagonal of a
directed adjacency matrix.

The known extremal instances have two realizations
(Proposition~\ref{shp:two}). Whether equality or near-equality holds in
either bound of \eqref{eq:gaps} for some degree vectors with large $M$ is not
known. How much larger the gaps are for non-extremal degree sequences, or
for typical sequences under a specified random model, is not quantified
here.

\appendix
\section{Three-row and graph-triple estimates}\label{app:section}
This appendix proves the triple estimate $H_T^2\succeq H_T$ used in
Proposition~\ref{asm:global}. The three-row variance inequality,
Lemma~\ref{app:rows}, is the conditional-projection inequality of Fu, Qin and
Wang~\cite[Theorem~5.1]{FQW26}, in the equivalent form \eqref{app:rowvariance}.
We reproduce the alternative proof of OpenAI~\cite[Section~3]{OAI26}. Its
fixed-block step, Lemma~\ref{app:fixed}, uses transposition eigenspaces, a
strategy that~\cite[Section~3]{OAI26} credits to the symmetric-group method
of Carlen, Lieb and Loss~\cite[Section~2]{CLL06}. Its count step,
Lemma~\ref{app:counts}, uses an eigenvector increment argument with the
adjacent-count coupling of Fu, Qin and Wang~\cite[Lemma~5.2]{FQW26}; the count
law corresponds to their Lemmas~5.2--5.4. The transfer to graph triples,
Lemma~\ref{app:triple}, uses the triple encoding of
OpenAI~\cite[Proposition~4.1]{OAI26}. All steps are included.

\subsection{A variance formulation}
For three coordinates on any finite probability space let $V_i$ be the
centered functions of coordinate $i$, and put the sum of marginal squared
norms on $V_1\oplus V_2\oplus V_3$. For
$A(g_1,g_2,g_3)=g_1+g_2+g_3$ the adjoint on centered joint functions is
$A^*f=(P_1f,P_2f,P_3f)$, where $P_i$ is conditional expectation onto that
coordinate. This follows by testing inner products. Therefore the assertions
\begin{equation}\label{app:rowvariance}
 \operatorname{Var}\Big(\sum_{i=1}^3g_i\Big)\le2\sum_{i=1}^3\operatorname{Var}(g_i),
 \qquad \sum_{i=1}^3P_i\preceq2I\text{ on }\mathbf 1^\perp
\end{equation}
are equivalent, since $AA^*=\sum_iP_i$ and both assertions say
$\lVert A\rVert^2\le2$.

\subsection{Fixed block sizes}
\begin{lemma}[Fixed block sizes]\label{app:fixed}
Take independent uniform ordered partitions of labeled sets $U,W$ into three
blocks of prescribed sizes, and let the row coordinates be $(U_i,W_i)$. Then
\eqref{app:rowvariance} holds for these coordinates.
\end{lemma}
\begin{proof}
Let $J$ be the sum of $I-T_\tau$ over label transpositions within either set.
These are self-adjoint involutions, and
$\langle f,Jf\rangle=\tfrac12\sum_\tau\lVert f-T_\tau f\rVert^2$.
The two permutation groups act transitively on the partition pairs, so
$\ker J=\mathbb R\mathbf 1$, including the degenerate one-state case.

Every transposition preserves each row-function space: it takes $U_i$ or
$W_i$ to its image under the known label permutation, without needing the
other blocks. Thus $J$ preserves these spaces. In finite dimension its
spectral projectors are polynomials in $J$ (interpolate on the distinct
eigenvalues), so they also preserve the row spaces. Center each $g_i$ and
project it into each positive eigenspace, writing $g_{i,\lambda}$.
At a given state any transposition changes at most two block tuples; hence
at most two of the three differences $g_{i,\lambda}-T_\tau g_{i,\lambda}$
are nonzero. The inequality $(a+b)^2\le2(a^2+b^2)$ gives
\[
 \lambda\Big\lVert\sum_i g_{i,\lambda}\Big\rVert^2
 \le2\lambda\sum_i\lVert g_{i,\lambda}\rVert^2.
\]
Cancel $\lambda>0$ and sum over orthogonal eigenspaces to obtain
\eqref{app:rowvariance}. The invariance of the row spaces is what permits
passing from the transposition energy estimate to this variance estimate.
\end{proof}

\subsection{Random block sizes}
\begin{lemma}[Random block sizes]\label{app:counts}
For nonnegative integers $m$ and $\alpha_1,\alpha_2,\alpha_3$ let
$w_i(t)=1/(t!(t+\alpha_i)!)$, and condition the product weights
$\prod_iw_i(x_i)$ on $x_1+x_2+x_3=m$. Then \eqref{app:rowvariance} holds for
the coordinates $x_i$.
\end{lemma}
\begin{proof}
All compositions of $m$ have positive probability. First consider two
coordinates of total $s$, with $p_s(y)\propto w_j(y)w_k(s-y)$,
$0\le y\le s$. On $0\le y<s$,
\[
 \frac{p_s(y)}{p_{s-1}(y)}
 =\frac{c_s}{(s-y)(s-y+\alpha_k)},\qquad c_s>0,
\]
which increases with $y$. Reweighting by an increasing positive likelihood
ratio increases expectations of increasing functions: for increasing $a,b$
and independent $Y,Y'$ with the old law,
\[
 2\operatorname{Cov}(a(Y),b(Y))
 =\mathbb E[(a(Y)-a(Y'))(b(Y)-b(Y'))]\ge0.
\]
The extra mass at the new largest value $s$ preserves this stochastic order,
so $Y_{s-1}\le_{\rm st}Y_s$. Apply the same argument to $Z_s=s-Y_s$ to obtain
$Y_s\le_{\rm st}Y_{s-1}+1$. Common-uniform quantiles simultaneously couple
these variables as $Y_{s-1}\le Y_s\le Y_{s-1}+1$. Thus increasing the total
by one increases exactly one of the two integer coordinates by one.

Use the addition operator $A$ above for the centered marginal spaces.
A nonzero eigenvector $b$ of $A^*A$ at eigenvalue $\lambda$ satisfies
\[
 \lambda b_i(t)=b_i(t)+\sum_{j\ne i}\mathbb E[b_j(x_j)\mid x_i=t],
 \qquad 0\le t\le m.
\]
If $m>0$, set $\eta=\max_{i,t<m}|b_i(t+1)-b_i(t)|>0$; otherwise centering
would make $b$ zero. Adjacent values of $x_i$ change the other total by one.
The coupling just proved changes exactly one of those coordinates by one,
so the difference of their summed conditional expectations has absolute
value at most $\eta$. Subtracting the eigenvector equations gives
\[
 |\lambda-1|\,|b_i(t+1)-b_i(t)|\le\eta.
\]
At a maximizing increment, $|\lambda-1|\le1$ and hence $\lambda\le2$.
Apply this to a largest eigenvalue of the positive $A^*A$ to prove
\eqref{app:rowvariance}. If $m=0$, every variance is zero.
\end{proof}

\subsection{All binary three-row margins}
\begin{lemma}[Three-row variance inequality]\label{app:rows}
Take the uniform binary matrices with any feasible three-row margins. Then
\eqref{app:rowvariance} holds for the three row coordinates.
\end{lemma}
\begin{proof}
Delete columns of sum zero or three and subtract the sum-three count from
the row totals. Let $U$ be the sum-one columns and $W$ the sum-two columns.
Write $U_i$ for the unique-one locations in row $i$, $W_i$ for the
unique-zero locations, and $r_i$ for the adjusted row total. Each row is
equivalent to $(U_i,W_i)$, and the whole matrix to the two ordered
partitions. With $k_i=|W_i|$ and $\delta_i=r_i-|W|$, one has
$|U_i|=k_i+\delta_i$. At fixed $k$ there are exactly
\[
 \frac{|U|!\,|W|!}{\prod_i k_i!(k_i+\delta_i)!}
\]
matrices, and the two uniform partition laws are independent. All these
partitions are feasible by their unique-one and unique-zero interpretation.

Set $l_i=\max(0,-\delta_i)$, $x_i=k_i-l_i$, $\alpha_i=|\delta_i|$, and
$m=|W|-\sum_i l_i$. Nonemptiness gives $m\ge0$. The feasible $x$ are all
nonnegative compositions of $m$; block upper bounds follow from their
nonnegative sizes and fixed totals. Also
$k_i!(k_i+\delta_i)!=x_i!(x_i+\alpha_i)!$, so the count law is the one of
Lemma~\ref{app:counts}. For a row function $f_i$, the conditional mean
$b_i(k_i)=\mathbb E[f_i\mid k]$ depends only on $k_i$, because its marginal
blocks are independent uniform subsets of sizes $k_i+\delta_i,k_i$.
Total variance, Lemma~\ref{app:fixed}, and then Lemma~\ref{app:counts} give
\begin{align*}
 \operatorname{Var}\Big(\sum_i f_i\Big)
 &=\mathbb E\operatorname{Var}\Big(\sum_i f_i\Bigm| k\Big)
  +\operatorname{Var}\Big(\sum_i b_i\Big)\\
 &\le2\sum_i\mathbb E\operatorname{Var}(f_i\mid k)+2\sum_i\operatorname{Var}(b_i)
 =2\sum_i\operatorname{Var}(f_i).
\end{align*}
This proves \eqref{app:rowvariance} for every feasible three-row margin.
\end{proof}

\subsection{Encoding graph triples}
Fix $T=\{i,j,k\}$, the graph outside $T$, and each outside vertex's number of
neighbors in $T$. These data define cells preserved by pair updates inside
$T$. The degree sum in $T$ minus the sum of outside counts determines twice
the internal edge count $e$. Encode the outside incidences by three rows,
and augment them as in Table~\ref{app:encoding}, where $s_i$ is the internal
degree of vertex $i$.

\begin{table}[ht]
\centering
\begin{tabular}{c l l l}
\toprule
$e$ & Internal graph & Extra column & Prescribed row sum\\
\midrule
0 & empty & none & $d_i$\\
1 & edge between marked vertices & $s_i$, sum 2 & $d_i$\\
2 & path centered at marked vertex & $s_i-1$, sum 1 & $d_i-1$\\
3 & triangle & none & $d_i-2$\\
\bottomrule
\end{tabular}
\caption{Augmentation of the three-row encoding of a graph triple by its
internal edge count $e$.}
\label{app:encoding}
\end{table}

The extra column is a separate label and uniquely reconstructs the internal
graph. Conversely every binary matrix of these margins reconstructs a graph
in the cell: the row sums impose the three degrees, and the column sums and
fixed outside graph impose all other degrees. There is no extra symmetry or
structural-zero restriction. This is a bijection preserving uniform measure.

An augmented row $S_i$ fixes the external neighbors of $i$, hence its
internal degree $s_i$ by subtraction from $d_i$; its extra entry adds no
restriction. Together with the cell data it determines the full pair-fiber
signature for $\{j,k\}$: all edges away from that pair, its internal edge
$jk$ with indicator $e-s_i$, the other outside pair counts obtained by
subtracting incidences to $i$, and the pair count $s_i$ of vertex $i$ itself.
Conversely the signature determines the augmented row. In particular $s_i=1$
fixes a count, not a choice of $j$ or $k$. Both conditionings average exactly
the same sets, so $E_{\{j,k\}}=P_i=\mathbb E[\,\cdot\mid S_i]$ on the cell.

\begin{lemma}[Graph triples]\label{app:triple}
For every vertex triple $T$, $H_T^2\succeq H_T$.
\end{lemma}
\begin{proof}
Lemma~\ref{app:rows} and the adjoint formulation \eqref{app:rowvariance}
give $\sum_iP_i\preceq2I$ on centered cell functions. Thus
$H_T=3I-\sum_iP_i\succeq I$ there, while $H_T$ kills cell constants.
Self-adjointness makes its eigenvalues zero or at least one, and hence
$H_T^2\succeq H_T$. Taking the orthogonal direct sum over cells proves the
claim on the entire graph space, including one-state cells.
\end{proof}

\section*{Acknowledgments}
The manuscript and its Lean verification were produced with AI assistance.
The manuscript openai-math\#131~\cite{OAI26} on which they build is
machine-checked but not peer reviewed. The verification artifacts, namely the Lean project and the scripts, are
available at \url{https://github.com/sevenevesai/SharpSwitchGap}.

\begin{thebibliography}{99}

\bibitem{Caputo08}
P.~Caputo,
\emph{On the spectral gap of the Kac walk and other binary collision processes},
ALEA. Latin American Journal of Probability and Mathematical Statistics
\textbf{4} (2008), 205--222.
\href{https://alea.impa.br/articles/v4/04-10.pdf}{Primary article}.

\bibitem{CLL06}
E.~Carlen, E.~H.~Lieb, and M.~Loss,
\emph{An Inequality of Hadamard Type for Permanents},
Methods and Applications of Analysis \textbf{13} (2006), no.~1, 1--18.
\href{https://arxiv.org/abs/math/0508096v1}{arXiv:math/0508096v1}.

\bibitem{CBS18}
C.~J.~Carstens, A.~Berger, and G.~Strona,
\emph{A unifying framework for fast randomization of ecological networks with fixed (node) degrees},
MethodsX \textbf{5} (2018), 773--780.\newline
\href{https://doi.org/10.1016/j.mex.2018.06.018}{doi:10.1016/j.mex.2018.06.018}.
Extended version: \href{https://arxiv.org/abs/1609.05137v3}{arXiv:1609.05137v3}.

\bibitem{CK18}
C.~J.~Carstens and P.~Kleer,
\emph{Speeding up Switch Markov Chains for Sampling Bipartite Graphs
with Given Degree Sequence},
Approximation, Randomization, and Combinatorial Optimization.
Algorithms and Techniques (APPROX/RANDOM 2018),
Leibniz International Proceedings in Informatics \textbf{116} (2018), 36:1--36:18.
\href{https://doi.org/10.4230/LIPIcs.APPROX-RANDOM.2018.36}{doi:10.4230/LIPIcs.APPROX-RANDOM.2018.36}.
Preprint: \emph{Comparing the Switch and Curveball Markov Chains for Sampling
Binary Matrices with Fixed Marginals},
\href{https://arxiv.org/abs/1709.07290v3}{arXiv:1709.07290v3}.

\bibitem{CDG07}
C.~Cooper, M.~Dyer, and C.~Greenhill,
\emph{Sampling regular graphs and a peer-to-peer network},
Combinatorics, Probability and Computing \textbf{16} (2007), no.~4, 557--593.\newline
\href{https://doi.org/10.1017/S0963548306007978}{doi:10.1017/S0963548306007978}.

\bibitem{CDG12}
C.~Cooper, M.~Dyer, and C.~Greenhill,
\emph{Corrigendum: Sampling regular graphs and a peer-to-peer network},
preprint (2012),
\href{https://arxiv.org/abs/1203.6111v1}{arXiv:1203.6111v1}.

\bibitem{DS87}
P.~Diaconis and M.~Shahshahani,
\emph{Time to reach stationarity in the Bernoulli--Laplace diffusion model},
SIAM Journal on Mathematical Analysis \textbf{18} (1987), no.~1, 208--218.
\href{https://doi.org/10.1137/0518016}{doi:10.1137/0518016}.

\bibitem{DGU14}
M.~Dyer, C.~Greenhill, and M.~Ullrich,
\emph{Structure and eigenvalues of heat-bath Markov chains},
Linear Algebra and its Applications \textbf{454} (2014), 57--71.
\href{https://doi.org/10.1016/j.laa.2014.04.018}{doi:10.1016/j.laa.2014.04.018}.

\bibitem{Unified22}
P.~L.~Erd\H{o}s, C.~Greenhill, T.~R.~Mezei, I.~Mikl\'os,
D.~Solt\'esz, and L.~Soukup,
\emph{The mixing time of switch Markov chains: a unified approach},
European Journal of Combinatorics \textbf{99} (2022), 103421.
\href{https://doi.org/10.1016/j.ejc.2021.103421}{doi:10.1016/j.ejc.2021.103421}.

\bibitem{Filmus16}
Y.~Filmus,
\emph{An Orthogonal Basis for Functions over a Slice of the Boolean Hypercube},
Electronic Journal of Combinatorics \textbf{23} (2016), no.~1, P1.23.
\href{https://doi.org/10.37236/4567}{doi:10.37236/4567}.

\bibitem{FQW26}
W.~Fu, Q.~Qin, and G.~Wang,
\emph{Spectral Gap for the Binary Fixed-Margin Swap Chain},
preprint (2026), \href{https://arxiv.org/abs/2606.22636v2}{arXiv:2606.22636v2}.

\bibitem{GS18}
C.~Greenhill and M.~Sfragara,
\emph{The switch Markov chain for sampling irregular graphs and digraphs},
Theoretical Computer Science \textbf{719} (2018), 1--20.
\href{https://doi.org/10.1016/j.tcs.2017.11.010}{doi:10.1016/j.tcs.2017.11.010}.

\bibitem{Havel55}
V.~Havel,
\emph{Pozn\'amka o existenci kone\v{c}n\'ych graf\r{u}}
[A remark on the existence of finite graphs],
\v{C}asopis pro p\v{e}stov\'an\'i matematiky
\textbf{80} (1955), no.~4, 477--480.
\href{https://doi.org/10.21136/CPM.1955.108220}{doi:10.21136/CPM.1955.108220}.

\bibitem{KTV99}
R.~Kannan, P.~Tetali, and S.~Vempala,
\emph{Simple Markov-chain algorithms for generating bipartite graphs and tournaments},
Random Structures \& Algorithms \textbf{14} (1999), no.~4, 293--308.
\href{https://doi.org/10.1002/(SICI)1098-2418(199907)14:4\%3C293::AID-RSA1\%3E3.0.CO;2-G}{doi:10.1002/(SICI)1098-2418(199907)14:4\textless293::AID-RSA1\textgreater3.0.CO;2-G}.

\bibitem{OAI26}
OpenAI,
\emph{Polynomial mixing of the switch chain for every graphical degree sequence},
OpenAI Math Release preprint
\href{https://github.com/openai/math/tree/adc7f1241b42e322a6451854ab7e4b4c146bf78a/preprints/Polynomial-Mixing-of-the-Switch-Chain-for-Every-Graphical-Degree-Sequence-September-25-2026/build}{OAI:Polynomial-Mixing-of-the-Switch-Chain-for-Every-Graphical-Degree-Sequence-September-25-2026},
2026.

\bibitem{TY23}
K.~Tikhomirov and P.~Youssef,
\emph{Sharp Poincar\'e and log-Sobolev inequalities for the switch chain on
regular bipartite graphs},
Probability Theory and Related Fields \textbf{185} (2023), 89--184.
\href{https://doi.org/10.1007/s00440-022-01172-7}{doi:10.1007/s00440-022-01172-7}.

\bibitem{Verhelst08}
N.~D.~Verhelst,
\emph{An efficient MCMC algorithm to sample binary matrices with fixed marginals},
Psychometrika \textbf{73} (2008), no.~4, 705--728.
\href{https://doi.org/10.1007/s11336-008-9062-3}{doi:10.1007/s11336-008-9062-3}.

\end{thebibliography}

\end{document}
